\chapter*{前言}

这本书从硬件本身开始。第一步不是急着给出一个大盒子，而是先把底层问题讲清楚：一根线上连续变化的电压怎样被稳定解释成 0 和 1；一个器件怎样让电流更容易沿某个方向通过；一个受控开关怎样让一个信号决定另一条电流路径是否导通。

沿着这条线继续向前，开关会组成逻辑门，逻辑门会组成加法器、比较器、选择器和译码器。组合逻辑能根据当前输入给出输出，却不能保存过去。为了让硬件一步一步推进，必须引入反馈、锁存器、触发器、时钟、复位、寄存器和状态机。

等状态边界和运算部件都讲清楚后，再进入 ALU、数据通路和控制器。最小 CPU 不在开头出现，而是在寄存器、ALU、数据通路和控制器都已经自然形成之后才出现。这样读者看到 CPU 时，它不是一个突然降临的大盒子，而是一组已经理解过的硬件部件按时钟边界组合起来。

最后再把这个最小计算核心接向主存、总线、设备控制器、中断控制器和启动硬件。读完整本书后，读者能把一台机器从电平、开关、门电路、组合逻辑、时序逻辑、运算部件、数据通路、控制器一路连到整机，而不是只记住一串名词。

本书采用接近 CSAPP（Computer Systems: A Programmer's Perspective）等经典系统教材的读法：每讲一个抽象，都同时给出机器层证据。读逻辑门时看真值表，也看开关路径；读组合逻辑时看函数，也看延迟和毛刺；读寄存器时看 bit，也看采样边界；读 CPU 时看指令，也看控制信号和数据通路 trace。章节末尾的“经典教材标注”不是额外习题，而是提醒读者用证据检查自己是否真的理解了硬件。
